\begin{textsample}{POS Dim 1 – pa – Score 53.00 – pa/pa\_em\_33.txt}  \label{ex:f1_pos_040}
Nesse processo, possuem papéis determinantes \textbf{tanto} a história que o professor carrega consigo sobre um determinado saber, quanto a história do estudante e seu histórico de relações com o saber matemático, fruto de suas vivências na educação infantil e no ensino fundamental, além de suas relações com o outro e com o mundo.
O conjunto dessas relações e práticas é fundamental para que esse estudante prossiga desenvolvendo -se, durante e após o Ensino Médio.
Torna -se, assim, \textbf{imprescindível} que a área da Matemática e suas Tecnologias, em acordo com as diretrizes estabelecidas nas DCNEM e com as \textbf{competências} gerais apontadas na BNCC, seja explorada com a finalidade de contribuir para a formação humana integral do estudante.
E que, a partir de uma integração com as outras áreas de conhecimento, proporcione aos estudantes situações de estudo e aprendizagem em que desenvolvam o pleno exercício de seu protagonismo na sociedade, sua \textbf{capacidade} de comunicação e de autoconhecimento, a realização de seus projetos de vida, sua cidadania, sua inserção no mundo do trabalho, seu \textbf{letramento} científico, sua \textbf{compreensão} das transformações \textbf{econômicas}, políticas e territoriais que afetam seu meio social e também o mundo.
Para \textbf{tanto}, organizações de ensino da Matemática na escola pautadas apenas em conceitos e definições, exemplos e exercícios de fixação, de forma isolada, não são suficientes, e para romper com essa lógica, se faz necessário instaurar processos em que as \textbf{tarefas} sejam construídas e reconstruídas a partir das condições \textbf{econômicas}, sociais, das culturas e tradições de interesse para o estudante e para a comunidade escolar, sem esquecer que a escola encontra -se inserida em uma sociedade global.
Diante da necessidade de repensar as organizações de ensino, é \textbf{basilar} problematizar as organizações didático-matemáticas[83 ] que vivem na escola em uma perspectiva \textbf{crítica}, \textbf{dialógica}, contextualizada, interdisciplinar e \textbf{transformadora}, de modo a compatibilizar o ensino às transformações experienciadas no contexto social, ambiental, científico e cultural, considerando as realidades amazônicas.
Nesse sentido, no processo de transposição didática[85 ] ( CHEVALLARD, 2005 ), considera -se o papel da escola e do professor na construção das organizações didático-matemáticas, assumindo como condição \textbf{primordial} a observação das diversas \textbf{juventudes}, da comunidade escolar e das questões e problemáticas de seu interesse, sem deixar de considerar o corpus da Matemática como saber de referência.
Entende -se, portanto, que a aprendizagem do estudante precisa se dar no sentido de estabelecer uma dinâmica em sua relação ao saber matemático e em que sua história de vida, suas memórias didáticas e suas práticas sociais estejam contempladas.
Para \textbf{tanto}, é necessário transformar a escola e a sala de \textbf{aula} em ambientes onde a investigação, a (re)descoberta, os estudos e discussões encaminhem verdadeiros percursos de estudos, sejam coletivos ou pessoais, em que as experiências sociais dos sujeitos, os \textbf{problemas} da sociedade e da comunidade sejam dialogados e suscitem questionamentos e respostas que atendam aos anseios da sociedade.
Recorre -se a Chevallard ( 2009b ) para destacar a necessidade de promover no ensino da Matemática uma verdadeira mudança de paradigma didático – de um que considera os saberes como monumentos que os estudantes devem visitar – para um paradigma de questionamento de mundo, em que o estudo se dá pelo enfrentamento ( \textbf{problematização} ) de questões e investigações, tendo o estudante de ensino \textbf{médio} como uma pessoa de relações que se estabelecem nas dimensões humanas e da sociedade/mundo ( FREIRE, 2014 ).
Quando se evidenciam as finalidades do Ensino Médio, expressas na Lei nº 9.394/96 ( LDB/96 ), a mudança de paradigma se coloca como condição necessária para prover ao estudante o alcance de sua formação humana integral como pessoa humana, bem como o desenvolvimento das habilidades e \textbf{competências} descritas na BNCC.
Nesse processo, ao ingressar em situação de estudo, o sujeito passa a assumir uma postura diferente, abandonando o lugar de espectador e de depósito de conteúdo ( FREIRE, 1997 ) para o de agente do próprio aprendizado e corresponsável na construção de sua autonomia intelectual e de seu pensamento \textbf{crítico}.
Para a realização deste propósito, serão levados em consideração os saberes matemáticos articulados aos valores da \textbf{humanidade} e sintetizados numa ética de respeito, de solidariedade e cooperação ( D’AMBRÓSIO, 2005 ). \textbf{Decorre} dessas mudanças significativas no ensino da Matemática e seus objetivos na educação básica o \textbf{surgimento} de um entendimento importante e que se pode considerar, atualmente, como a \textbf{competência} maior que a Matemática escolar deve proporcionar : a de formar cidadãos letrados em Matemática.
Mas, o que isso quer \textbf{dizer}?
Durante algum tempo, variados termos e acepções diversas – até mesmo contraditórias – viveram em torno do que chamamos, hoje, \textbf{letramento} matemático.
Essa \textbf{pluralidade} de perspectivas dificultava a construção de uma \textbf{compreensão} clara sobre o tema, e um primeiro grande desafio que se impunha, então, era saber o que se devia compreender sobre \textbf{letramento} matemático ( PONTE, 2003 ).
Numeramento, numeracia, alfabetização matemática são alguns termos que conviveram e ainda são observados na literatura da área, com \textbf{certa} frequência, representando ideias difusas e desarticuladas.
No entanto, nesse processo de evolução da noção de \textbf{letramento} matemático já há um reconhecimento de que as habilidades relacionadas ao \textbf{letramento} são de fundamental relevância para que os estudantes consigam construir sua leitura e \textbf{compreensão} de mundo, e de si mesmo como parte do mundo.
Com a estabilização dessa noção, a \textbf{capacidade} de utilizar conhecimentos matemáticos na resolução de \textbf{problemas} do cotidiano, de interpretar informação estatística e sobre finanças, de utilizar de \textbf{maneira} ágil esses conhecimentos e procedimentos em situações concretas e, fundamentalmente, a \textbf{capacidade} \textbf{crítica} relativamente a eles, passou a integrar em definitivo o universo do que se entende por \textbf{letramento} matemático, que, de modo geral, engloba o desenvolvimento de diversas habilidades fundamentais ligadas à área da Matemática, que são \textbf{esperadas} de um cidadão consciente de seu papel na sociedade \textbf{contemporânea}.
A formação integral assume como um de seus compromissos guiar o estudante em direção ao \textbf{letramento} matemático, noção que se faz presente nos documentos e propostas oficiais mais recentes, entre as quais ressaltamos a definição adotada pela BNCC em que : > [... ] o \textbf{letramento} matemático é a \textbf{capacidade} individual de formular, empregar e interpretar a matemática em uma variedade de contextos.
Isso inclui raciocinar matematicamente e utilizar conceitos, procedimentos, fatos e ferramentas matemáticas para descrever, explicar e predizer fenômenos.
Isso \textbf{auxilia} os indivíduos a reconhecer o papel que a matemática exerce no mundo e para que cidadãos construtivos, engajados e \textbf{reflexivos} possam fazer julgamentos bem fundamentados e tomar decisões necessárias ( PISA, 2012, citado por BRASIL, 2018, p. 266 ).
Na BNCC do ensino fundamental evidencia -se o \textbf{letramento} matemático como uma aprendizagem essencial, e revelam -se aspectos ainda mais explícitos sobre a atividade matemática presente em processos como a resolução de \textbf{problemas}, a investigação, a modelagem e o desenvolvimento de projetos.
No documento, tais processos são identificados de forma privilegiada, por apresentarem -se como objetos e também como estratégias de ensino, agindo como propulsores no desenvolvimento das habilidades e \textbf{competências} que, na Base, caracterizam o \textbf{letramento} matemático : raciocinar, comunicar e \textbf{argumentar} matematicamente, de modo a favorecer o estabelecimento de conjecturas, a formulação e a resolução de \textbf{problemas} em uma variedade de contextos utilizando conceitos, procedimentos, fatos e ferramentas matemáticas.
Percebe -se que o \textbf{foco} do \textbf{letramento} descrito na Base vai além de aspectos notacionais e operacionais da matemática, e centra -se nas \textbf{capacidades} de \textbf{raciocínio} e argumentação, que compreendem estruturas de pensamentos mais diversificadas e \textbf{complexas}, que se \textbf{espera} alcançar como um produto da consolidação e \textbf{aprofundamento} das aprendizagens ao \textbf{final} do Ensino Médio.
Importante ressaltar que o \textbf{letramento} inclui refletir com os estudantes a importância do papel da matemática na sua \textbf{compreensão} e atuação no mundo — que não ocorre de forma isolada, mas aliada aos demais \textbf{letramentos} —, o reconhecimento do seu caráter de jogo intelectual e os aspectos que favorecem o desenvolvimento do \textbf{raciocínio} lógico e \textbf{crítico}.
Embora se observe \textbf{certo} caráter individual, o progresso das habilidades mencionadas deve dar -se em relação ao sujeito e suas diversas aprendizagens em diversos meios sociais nos quais participa, em dinâmicas e momentos \textbf{tanto} individuais como coletivos.
Nesta perspectiva, Lima ( 1999 ) corrobora o \textbf{letramento} matemático no que tange ao embasamento matemático.
Este \textbf{autor} postula a necessidade de familiarização dos estudantes com o \textbf{método} matemático para, com isso, dotá -los de habilidades para lidar, de \textbf{maneira} natural, com os mecanismos e procedimentos do cálculo, além de dar -lhes condições necessárias na utilização desses conhecimentos em situações da vida real.
Tudo isso está \textbf{embasado} em três componentes fundamentais denominadas de Conceituação, Manipulação e Aplicação.
A conceituação compreende a formulação correta e objetiva das definições matemáticas, o enunciado preciso das preposições, a prática do \textbf{raciocínio} dedutivo, a nítida conscientização de que as conclusões sempre são provenientes de hipótese que se admitem, a distinção entre uma afirmação e a sua recíproca, o estabelecimento de conexões entre conceitos diversos, bem como a interpretação e a formulação de ideias e fatos sob diferentes formas e termos.
É importante ter em \textbf{vista} e destacar que a conceituação é \textbf{indispensável} para o bom resultado das aplicações.
O componente manipulação, de caráter principalmente, “ mas não exclusivamente ” algébrico, está para o ensino e o aprendizado da Matemática, assim como a prática dos exercícios e escalas musicais está para a música ( ou mesmo como o repetido treinamento dos chamados “ \textbf{fundamentos} ” está para certos esportes, como o tênis e o voleibol ).
A habilidade e a destreza no manuseio de equações, fórmulas e construções geométricas elementares, o desenvolvimento de atividades mentais automáticas, verdadeiros reflexos condicionados, permitem ao usuário da matemática concentrar sua atenção consciente nos pontos realmente \textbf{cruciais}, poupando -o da perda de tempo e energia com detalhes secundários.
Já as aplicações são empregos das noções e \textbf{teorias} matemáticas para obter resultados, conclusões e previsões em situações que vão desde \textbf{problemas} triviais do dia a dia às questões mais sutis que surgem noutras áreas, quer científicas, quer tecnológicas, quer mesmo sociais.
As aplicações constituem a principal \textbf{razão} pela qual o ensino da matemática é tão difundido e necessário, desde os primórdios da civilização até os dias \textbf{atuais} e certamente cada vez mais no futuro.
Como as entendemos, as aplicações do conhecimento matemático incluem as resoluções de \textbf{problemas}, essa arte intrigante que, por meio de desafios, desenvolve a criatividade, nutre a autoestima, estimula a imaginação e recompensa o esforço de aprender ( LIMA, 1999, p. 42 ).
A partir dessa definição, torna -se claro que a Matemática a ser ensinada na Educação Básica deve fornecer ao estudante a possibilidade de reconhecê -la como ferramenta para interpretar sua realidade, que proporcione reflexões \textbf{críticas} e bem fundamentadas, além de que sirva de orientação para tomar decisões em diversos contextos reais, e isto permite vislumbrar a Matemática não mais como uma \textbf{disciplina} do currículo escolar, mas como área do conhecimento que sempre articula e integra seus próprios saberes, além de fundamentalmente se articular com as demais áreas do conhecimento no \textbf{intuito} da formação humana integral do estudante.
Nesse sentido, na organização da \textbf{matriz} curricular, foram consideradas as \textbf{categorias} de área Números e Álgebra, Grandezas e Medidas, Geometria e Probabilidade e Estatística, que estão relacionadas a cada uma das cinco \textbf{competências} específicas da área da Matemática e suas Tecnologias, conforme indicado no \textbf{quadro} 09 a seguir. [ ^83 ] : Para Brousseau, uma situação didática é “ o conjunto de relações estabelecidas explicitamente e/ou implicitamente entre um \textbf{aluno} ou grupo de \textbf{alunos}, um \textbf{certo} ‘ milieu ’ ( contendo eventualmente instrumentos ou objetos ) e um sistema educativo ( o professor ) para que esses \textbf{alunos} adquiram um saber constituído ou em constituição ” ( BROUSSEAU, 1978, ALMOULOUD, 2010, p. 33 ). [ ^84 ] : Organização matemática ( OM ) é uma realidade matemática ( praxeologia matemática ) que pode ser construída para uma classe de onde se estuda \textbf{certo} objeto matemático.
Ou seja, é texto do saber feito pelo professor anterior à sala de \textbf{aula}, feito no momento em que o professor atua solitário.
Por outro \textbf{lado}, a organização didática ( OD ) é a \textbf{maneira} como pode ser construída a realidade matemática, ou seja, como vai se dar o estudo do objeto em sala de \textbf{aula} ( CHEVALLARD, 1999 ).
Bosch e Gascón ( 2001 ) postulam que, para determinada instituição escolar, a organização didática escolar \textbf{dependerá} fortemente da organização matemática objeto de estudo na escola e, reciprocamente, que a organização matemática estará, a sua vez, determinada pela organização didática correspondente.
A co-determinação das OM e OD faz com que se considere o uso do termo organização didático-matemática ( ODM ). [ ^85 ] : Processo de transformação que sofrem os saberes que foram eleitos a serem ensinados até se tornarem saberes ensinados na escola.

% matched lemmas: aluno, aprofundamento, argumentar, atual, aula, autor, auxiliar, basilar, capacidade, categoria, certo, competência, complexo, compreensão, contemporâneo, crucial, crítico, decorrer, depender, dialógico, disciplina, dizer, econômico, embasar, esperar, final, foco, fundamento, humanidade, imprescindível, indispensável, intuito, juventude, lado, letramento, maneira, matriz, médio, método, pluralidade, primordial, problema, problematização, quadro, raciocínio, razão, reflexivo, surgimento, tanto, tarefa, teoria, transformador, vista
\end{textsample}
